% =============================================================================
%  Manual del Aula Virtual de Física
%  Compilar:  pdflatex manual_aula_virtual.tex   (dos veces, para el índice)
%  Paquetes estándar de TeX Live (texlive-latex-extra, texlive-lang-spanish) o Overleaf.
% =============================================================================
\documentclass[11pt,a4paper,oneside]{book}

\usepackage[utf8]{inputenc}
\usepackage[T1]{fontenc}
\usepackage{lmodern}
\usepackage[spanish,es-noshorthands,es-tabla]{babel}
\usepackage[margin=2.5cm,headheight=14pt]{geometry}
\usepackage{amssymb,amsmath}
\usepackage{xcolor}
\usepackage{booktabs,tabularx,array,longtable}
\usepackage{enumitem}
\usepackage{fancyhdr}
\usepackage{titlesec}
\usepackage{tcolorbox}
\tcbuselibrary{breakable}
\usepackage[hidelinks,pdftitle={Manual del Aula Virtual de Física},pdfauthor={Aula Virtual de Física}]{hyperref}

\definecolor{primario}{HTML}{14284B}
\definecolor{acento}{HTML}{C2410C}
\definecolor{verde}{HTML}{047857}
\definecolor{gris}{HTML}{4B5563}

\titleformat{\chapter}[display]{\normalfont\huge\bfseries\color{primario}}{\chaptertitlename\ \thechapter}{12pt}{\Huge}
\titleformat{\section}{\normalfont\Large\bfseries\color{primario}}{\thesection}{1em}{}
\titleformat{\subsection}{\normalfont\large\bfseries\color{primario}}{\thesubsection}{1em}{}

\pagestyle{fancy}
\fancyhf{}
\fancyhead[L]{\small\color{gris}Aula Virtual de Física}
\fancyhead[R]{\small\color{gris}\leftmark}
\fancyfoot[C]{\thepage}
\renewcommand{\headrulewidth}{0.4pt}

\setlist{itemsep=2pt,topsep=4pt}
\setlength{\parskip}{4pt}

% Recuadros
\newtcolorbox{consejo}[1][Consejo]{colback=verde!5,colframe=verde,fonttitle=\bfseries,title=#1,breakable}
\newtcolorbox{importante}[1][Importante]{colback=acento!5,colframe=acento,fonttitle=\bfseries,title=#1,breakable}
\newtcolorbox{nota}[1][Nota]{colback=primario!4,colframe=primario,fonttitle=\bfseries,title=#1,breakable}

% Elementos de la interfaz
\newcommand{\boton}[1]{\textsf{\textbf{[\,#1\,]}}}
\newcommand{\menu}[1]{\textsf{\textbf{#1}}}
\newcommand{\ruta}[1]{\textsf{#1}}

\begin{document}

% ----------------------------------------------------------------------------
\begin{titlepage}
  \centering
  \vspace*{3cm}
  {\color{primario}\rule{\linewidth}{1.5pt}\par\vspace{0.6cm}}
  {\Huge\bfseries\color{primario} Aula Virtual de Física\par}
  \vspace{0.4cm}
  {\LARGE Manual de uso\par}
  {\color{primario}\vspace{0.6cm}\rule{\linewidth}{1.5pt}\par}
  \vspace{1.5cm}
  {\large Para estudiantes y profesorado\par}
  \vspace{0.8cm}
  {\normalsize Evaluaciones de una pregunta a la vez · Prácticas con aprendizaje asistido\\
  Seguimiento académico · Centro de Datos\par}
  \vfill
  {\normalsize Versión del 1 de noviembre de 2026\par}
\end{titlepage}

\tableofcontents

% ============================================================================
\chapter{Introducción}

\section{Qué es el Aula Virtual}
El Aula Virtual de Física es la plataforma del curso: teoría por lecciones, prácticas, evaluaciones, asistencia,
entregas en PDF, calificaciones y comunicación con el profesor. Funciona en el navegador del teléfono o de la
computadora; no hace falta instalar nada (también puede añadirse a la pantalla de inicio del teléfono).

Hay tres perfiles:
\begin{itemize}
  \item \textbf{Estudiante}: aprende, practica, presenta evaluaciones y ve sus resultados.
  \item \textbf{Profesor}: crea actividades, sigue a sus estudiantes, revisa y califica.
  \item \textbf{Administrador}: además gestiona las cuentas y la puesta en producción.
\end{itemize}

\section{Entrar}
Abra la dirección del aula y pulse \boton{Ingresar}. Use el identificador y la contraseña que le dio el profesor.
La primera vez se le pedirá cambiar la contraseña. Si la olvida, pida al profesor que la restablezca.

\section{Moverse por el aula}
\begin{itemize}
  \item En el \textbf{teléfono}, la barra inferior lleva a las secciones principales.
  \item En la \textbf{computadora}, la barra superior hace lo mismo.
  \item El menú de su nombre (arriba a la derecha) contiene su cuenta, la \menu{Ayuda}, el Laboratorio virtual y
        \boton{Cerrar sesión}.
  \item El botón del sol o la luna alterna entre el modo claro y el oscuro.
\end{itemize}

\begin{consejo}[Ayuda dentro del aula]
La página \menu{Ayuda} (menú de su nombre) resume este manual y se adapta a su perfil.
\end{consejo}

% ============================================================================
\part{Para estudiantes}

\chapter{Navegación del estudiante}

\begin{center}
\begin{tabularx}{\linewidth}{>{\bfseries}l X}
\toprule
Sección & Para qué sirve \\
\midrule
Inicio & Lo pendiente, sus cursos, avisos y el acceso para continuar donde lo dejó. Muestra por separado las
         \emph{evaluaciones} y las \emph{prácticas} pendientes. \\
Aprender & El temario: teoría por lecciones y la práctica libre de cada lección. \\
Practicar & Las prácticas asignadas por el profesor (se pueden repetir y tienen retroalimentación). \\
Evaluaciones & Exámenes y quizzes calificables (se entregan una vez). \\
Asistencia & Marcar su presencia en clase y ver su registro. \\
Mi cuenta & Calificaciones, progreso por asignatura, entregas en PDF, laboratorio, perfil y contraseña. \\
\bottomrule
\end{tabularx}
\end{center}

Cada actividad muestra su estado con una etiqueta y una sola acción:
\textbf{Disponible} (\boton{Empezar} o \boton{Practicar}), \textbf{En curso} (\boton{Continuar}),
\textbf{Próximamente} (aún no abre), \textbf{Entregada} (\boton{Ver}) o \textbf{Cerrada}. La fecha de cierre aparece en
rojo cuando faltan menos de 48 horas.

\chapter{Practicar}

\section{Práctica tradicional}
Todos los ejercicios aparecen en una página. Arriba hay un \textbf{mapa de ejercicios}: indica cuántos lleva
respondidos y permite saltar a cualquiera tocando su número (verde = respondido; con marca de verificación = correcto).

En cada ejercicio puede \boton{Comprobar} su respuesta (según lo que permita el profesor), pedir pistas y, en algunas
prácticas, ver la solución. Sus respuestas se guardan solas.

\section{Práctica con aprendizaje asistido}
En estas prácticas puede usar inteligencia artificial, libros, apuntes o la ayuda de otra persona para
\emph{aprender}. El aula le guía por cuatro momentos en cada ejercicio, uno a la vez, con una barra de progreso:

\begin{enumerate}[label=\textbf{\arabic*.}]
  \item \textbf{Intentar.} Resuelva el ejercicio por su cuenta. Puede escribir sus operaciones o marcar
        \emph{«No sé resolverlo todavía»}: no necesita acertar para continuar. Al pulsar
        \boton{Guardar mi primer intento} se le pregunta cómo lo hizo:
        \begin{itemize}
          \item Lo resolví por mí mismo, sin utilizar IA ni ayuda externa.
          \item Consulté alguna ayuda antes de responder.
          \item No conseguí resolverlo por mi cuenta.
        \end{itemize}
        Pulse \boton{Confirmar y guardar}. Su primera respuesta queda guardada tal cual y ya no se puede cambiar.
  \item \textbf{Aprender.} Ahora puede consultar cualquier ayuda permitida. Verá \emph{Tu respuesta inicial} (fija) y
        \emph{Tu respuesta después de estudiar} (editable). Corrija lo necesario y pulse \boton{Guardar mi corrección}:
        recibirá retroalimentación. Hay un número limitado de comprobaciones (normalmente dos).
  \item \textbf{Explicar.} Con sus palabras, responda a dos o tres preguntas breves; cambian según su caso (si se
        equivocó, si acertó o si no sabía). Marque las ayudas que usó. Si quiere, añada una referencia (fuente,
        enlace, concepto o idea que le ayudó). No hace falta copiar conversaciones privadas. Pulse
        \boton{Guardar mi explicación}. Después verá la solución.
  \item \textbf{Comprobar} (si el profesor lo activó). Un ejercicio \emph{distinto} con el mismo concepto. Tiene una
        sola oportunidad: \boton{Comprobar mi aprendizaje}.
\end{enumerate}

Al terminar todos los ejercicios, pulse \boton{Finalizar la práctica}.

\begin{importante}[Sobre la declaración de ayuda]
Declarar que usó ayuda \textbf{no le quita puntos}, y declarar que no la usó \textbf{no le da puntos extra}. Es una
información voluntaria para que el profesor entienda mejor cómo aprende el grupo; no es una verificación técnica.
\end{importante}

\begin{nota}[Sin conexión]
Si se va la internet, su borrador queda guardado en el teléfono y el aula no avanza de paso hasta que el servidor
confirme. Cuando vuelva la red, pulse otra vez el botón: no se guardará dos veces.
\end{nota}

\chapter{Evaluaciones}

Hay dos tipos de evaluación:
\begin{itemize}
  \item \textbf{Clásica}: todas las preguntas en una página; se entrega con \boton{Entregar}.
  \item \textbf{De una pregunta a la vez}: la que se describe en este capítulo.
\end{itemize}

\section{Antes de empezar}
Al abrir la evaluación verá el recuadro \textbf{«Antes de empezar»} con las reglas: duración, si se puede volver
atrás, si hay pausas y cuántas. Si hay límite de tiempo, el botón dice \boton{Comenzar (empieza el reloj)}.

\begin{importante}
El reloj lo lleva el \textbf{servidor}: sigue corriendo aunque cierre la página o apague el teléfono (salvo durante
una pausa permitida que detenga el reloj). Cambiar la hora del dispositivo no da ni quita tiempo.
\end{importante}

\section{Durante la evaluación}
\begin{itemize}
  \item Arriba: preguntas respondidas, la \textbf{cuenta atrás} y el estado del guardado
        (\emph{Guardado}, \emph{Guardando\dots}, \emph{Sin conexión: guardado en este dispositivo}).
  \item Debajo: los números de las preguntas (verde = respondida; con bandera = marcada para revisar). Tocando un
        número va a esa pregunta, si la navegación lo permite.
  \item Abajo: \boton{Anterior}, \boton{Marcar} y \boton{Siguiente}. En la última pregunta,
        \boton{Revisar y entregar}.
  \item Durante la evaluación se oculta la barra de navegación del aula (modo concentración), para no salir por
        accidente.
\end{itemize}

\textbf{Navegación.} En la navegación \emph{libre} puede volver a cualquier pregunta. En la navegación
\emph{solo hacia adelante}, una pregunta que deja atrás no se puede volver a abrir ni cambiar.

\textbf{Pausas.} Si el profesor las permite, verá \boton{Pausar}. Durante la pausa las preguntas se ocultan y no se
puede responder. Pulse \boton{Reanudar} para seguir.

\textbf{Prórrogas.} Si el profesor le concede más tiempo, aparecerá el aviso \emph{«El profesor amplió su tiempo»} y
la cuenta atrás aumentará.

\section{Entregar}
La pantalla \textbf{Revisión antes de entregar} muestra qué preguntas faltan y cuáles marcó. Pulse
\boton{Entregar el examen} y confirme. Si el tiempo se acaba, el aula entrega automáticamente lo que haya guardado
(antes consulta al servidor por si hubo una prórroga).

\section{Si se va la internet o se cierra la página}
Siga respondiendo: lo que escriba se guarda en el dispositivo y se envía solo al volver la red. Si se cerró la página,
vuelva a entrar en la evaluación: continúa en la misma pregunta con sus respuestas.

\chapter{Escribir respuestas}

\section{Respuestas numéricas}
Escriba el número y la unidad, por ejemplo \texttt{9.8 m/s\textasciicircum2}, \texttt{1.2e3 N} o
\texttt{10290 J}. Puede usar coma o punto decimal y las teclas rápidas bajo el campo.

\begin{itemize}
  \item Respete las mayúsculas de los símbolos: J, N, W, Pa, K.
  \item Si escribe \texttt{10290j}, normalmente se entenderá como julios y se le avisará de que el símbolo
        correcto es «J» (en algunas preguntas el profesor puede exigir el símbolo exacto).
  \item Si la escritura es ambigua (por ejemplo \texttt{mj}, que podría ser mJ o MJ), se le pedirá aclararla.
  \item Una unidad de otra magnitud (por ejemplo newtons en una pregunta de energía) es un error de concepto, no de
        escritura.
\end{itemize}

\section{Otros formatos}
Selección (una o varias opciones), verdadero o falso (a veces con justificación), texto corto y desarrollo
(lo revisa el profesor).

\chapter{Resultados, asistencia y entregas}
\begin{itemize}
  \item \textbf{Calificaciones} (\menu{Mi cuenta}): solo aparecen cuando el profesor las publica.
  \item \textbf{Mi progreso}: dominio por concepto en cada asignatura, calculado con sus respuestas verificables.
  \item \textbf{Asistencia}: cuando el profesor abre una sesión, márquela desde el aula durante la clase.
  \item \textbf{Entregas en PDF}: suba sus trabajos al buzón de la sección antes del cierre.
\end{itemize}

% ============================================================================
\part{Para el profesorado}

\chapter{Panel docente}

El panel muestra arriba lo de cada día:
\begin{itemize}
  \item \textbf{Crear una actividad}: un botón por sección.
  \item \textbf{Respuestas por revisar}: desarrollos y respuestas que el motor no puede calificar solo.
  \item \textbf{Pueden necesitar apoyo}: estudiantes con avisos de seguimiento, por sección.
\end{itemize}
Debajo: cifras generales, \textbf{Próximos cierres} de todas las secciones, sus secciones con accesos directos
(\boton{+ Actividad}, \boton{Estudiantes}, \boton{Notas}) y las herramientas: contenido, banco de ejercicios,
comunicación en vivo, Orbe Control, analítica de tiempo, vista de estudiante y Laboratorio LIEIA.

\section{Secciones}
Cada sección tiene pestañas: Estudiantes y grupos, Asistencia, Avisos, Actividades, \textbf{Seguimiento}, Entregas en
PDF, Calificaciones, Revisión manual y Estadísticas. Los estudiantes se matriculan por su identificador o con el
código de la sección.

\chapter{Crear actividades}

\section{El asistente en cuatro pasos}
\ruta{Panel docente $\rightarrow$ + Crear actividad}, o en la sección, pestaña Actividades.

\begin{enumerate}[label=\textbf{\arabic*.}]
  \item \textbf{Tipo}: elija una plantilla.
  \item \textbf{Preguntas}: márquelas en el banco (filtros por lección, dificultad y búsqueda) o déjelo para después.
  \item \textbf{Reglas}: título, fechas y, en evaluaciones, duración y navegación. El resto está en
        \emph{Opciones avanzadas}.
  \item \textbf{Revisar}: un resumen en lenguaje llano; pulse \boton{Crear y publicar} o \boton{Guardar borrador}.
\end{enumerate}

La actividad se crea oculta, se completa y solo al final se publica, para que nadie la empiece sin preguntas.

\begin{center}
\begin{tabularx}{\linewidth}{>{\bfseries}l X}
\toprule
Plantilla & Valores predeterminados \\
\midrule
Práctica libre & Retroalimentación al comprobar, pistas, tres comprobaciones por ejercicio, intentos ilimitados. \\
Práctica con aprendizaje asistido & Intentar $\rightarrow$ aprender $\rightarrow$ explicar ($\rightarrow$ comprobar);
                                     sin solución inmediata ni pistas; dos comprobaciones de la corrección; cierre en 7 días. \\
Tarea con fecha & Un intento; soluciones al cerrar; cierre en 7 días. \\
Quiz rápido & Una pregunta a la vez, 15 minutos, un intento, sin pausas. \\
Examen parcial & Una pregunta a la vez, 60 minutos, navegación libre, una pausa de 5 minutos que detiene el reloj. \\
Examen estricto & Una pregunta a la vez, 45 minutos, solo hacia adelante, sin pausas. \\
En blanco & Práctica o evaluación sin valores sugeridos. \\
\bottomrule
\end{tabularx}
\end{center}

\begin{consejo}
Antes de publicar una evaluación, ábrala con \boton{Vista previa como estudiante}.
\end{consejo}

\chapter{Evaluaciones de una pregunta a la vez}

\section{Modo de examen}
\ruta{Actividad $\rightarrow$ Configuración $\rightarrow$ Modo de examen}.

\begin{center}
\begin{tabularx}{\linewidth}{>{\bfseries}l X}
\toprule
Opción & Significado \\
\midrule
Presentación & \emph{Una pregunta a la vez} (reloj oficial, tiempo por pregunta, recuperación tras desconexiones) o
               \emph{Todas en una página} (clásico). \\
Navegación & \emph{Libre} (puede volver y marcar) o \emph{Solo hacia adelante} (lo impone el servidor). \\
Duración & Minutos desde que el estudiante empieza; nunca pasa del cierre de la actividad. \\
Pausas & Número máximo y minutos en total; las preguntas se ocultan durante la pausa. Si «detienen el reloj», el
         límite se desplaza por la pausa usada, hasta su máximo. \\
\bottomrule
\end{tabularx}
\end{center}

Los intentos ya empezados conservan la configuración con la que empezaron.

\section{Durante el examen: seguimiento en vivo}
\ruta{Actividad $\rightarrow$ Resultados $\rightarrow$ En vivo}. Se actualiza cada 10 segundos y muestra, por
estudiante: \emph{Respondiendo}, \emph{En pausa}, \emph{Sin conexión} (más de 45 segundos sin contacto) o
\emph{Tiempo vencido}; la pregunta en que va, cuántas lleva respondidas y el tiempo restante. No muestra respuestas.

\section{Prórrogas individuales}
En la fila del estudiante pulse \boton{+5 min}. Puede escribir un motivo (conexión caída, adaptación, incidencia); queda
registrado quién la dio, cuánto y cuándo. La prórroga nunca pasa del cierre de la actividad. El estudiante ve el aviso
en menos de 20 segundos.

\section{Después: tiempo por pregunta}
En la revisión de cada intento aparece el \textbf{tiempo por pregunta}, separado en trabajo, inactividad, pausa,
página oculta, desconexión y desconocido, con las veces que volvió a cada pregunta. Es una descripción del proceso,
no un indicador de conducta: una pausa o una página oculta tienen muchas explicaciones legítimas.

\chapter{Prácticas con aprendizaje asistido}

\section{Propósito}
El estudiante puede aprovechar la inteligencia artificial para aprender, y usted puede distinguir:
\begin{enumerate}
  \item lo que el estudiante \emph{declara} poder resolver solo (primer intento);
  \item lo que consigue después de estudiar con ayuda (corrección);
  \item qué errores identifica y qué dice haber aprendido (explicación);
  \item lo que demuestra después en un ejercicio diferente (transferencia y verificación posterior).
\end{enumerate}

\section{Crear y configurar}
\begin{enumerate}
  \item \ruta{+ Crear actividad $\rightarrow$ Práctica con aprendizaje asistido}. Si quiere el paso 4, marque
        \emph{Añadir el paso «Demuestra lo que aprendiste»}.
  \item \ruta{Actividad $\rightarrow$ Configuración $\rightarrow$ Aprendizaje asistido}:
    \begin{itemize}
      \item \textbf{Ejercicio de transferencia} de cada pregunta: un ejercicio validado del banco con el mismo concepto.
            Para comprobar comprensión, que no sea solo un cambio de números. Sin ejercicio asignado, esa pregunta no
            tiene paso 4.
      \item \textbf{Opciones avanzadas}: la ponderación y la verificación posterior (en otra sesión, en una fecha, en
            clase o en un examen). Si la vincula a una actividad, verá la \emph{retención} de cada estudiante.
    \end{itemize}
\end{enumerate}
Una práctica tradicional también puede convertirse desde esa tarjeta con \boton{Activar el aprendizaje asistido}.

\section{Calificación}
\begin{center}
\begin{tabular}{l r l}
\toprule
Fase & Peso & Cómo cuenta \\
\midrule
Intento inicial & 20\,\% & Por presentarlo (acertar no es necesario) \\
Corrección & 30\,\% & Resultado de la última versión \\
Explicación & 20\,\% & Presentada (al menos una frase); usted puede ajustar la calidad \\
Transferencia & 30\,\% & Resultado del ejercicio nuevo \\
\bottomrule
\end{tabular}
\end{center}

\begin{itemize}
  \item Los pesos son una \textbf{propuesta editable}, no una escala científica validada.
  \item Sin transferencia, los pesos activos se reparten proporcionalmente (28,6\,\% / 42,9\,\% / 28,6\,\%); la tarjeta
        de configuración lo muestra.
  \item \textbf{La declaración de uso de IA no cambia la nota.}
  \item La nota se integra en el libro de calificaciones como cualquier práctica, y puede corregirla con la
        calificación manual de la revisión.
\end{itemize}

\section{Proceso de aprendizaje}
\ruta{Actividad $\rightarrow$ Resultados $\rightarrow$ Proceso de aprendizaje}.

\textbf{Resumen de la sección}: primeros intentos presentados, porcentaje que declaró haber usado ayuda,
\emph{rendimiento inicial declarado}, \emph{rendimiento posterior a la asistencia}, explicaciones presentadas,
transferencia correcta, ayudas más usadas y ejercicios que más correcciones requirieron.

\textbf{Por estudiante} (una tabla por ejercicio): primer intento (con su resultado y la declaración), ayudas, respuesta
final, qué cambió (valor, unidad, opción, procedimiento), explicación (desplegable, con sus referencias), transferencia
y retención.

\begin{importante}[Cómo interpretar estos datos]
\begin{itemize}
  \item «Rendimiento inicial \emph{declarado}» no es «dominio autónomo verificado»: una práctica en casa no prueba
        trabajo autónomo.
  \item Una explicación correcta o bien redactada no prueba quién la escribió, y una que «parece de IA» no prueba nada.
  \item Ninguna cifra indica quién hizo trampa. Use la información para acompañar el aprendizaje.
\end{itemize}
\end{importante}

\section{Lo que garantiza el sistema}
\begin{itemize}
  \item \textbf{El primer intento no se puede cambiar}, ni siquiera directamente en la base de datos. Las
        correcciones y explicaciones son versiones nuevas.
  \item \textbf{No se puede tantear la respuesta antes del primer intento}: solo se comprueba que la respuesta se
        pueda leer, con un mensaje genérico; las pistas, la comprobación y la solución clásicas están cerradas.
  \item \textbf{Las fases van en orden} y cada estudiante solo accede a su propio trabajo.
  \item \textbf{Reintentos y dos pestañas}: un reenvío no se duplica; un segundo primer intento se rechaza.
\end{itemize}

\chapter{Revisar, calificar y publicar}
\begin{enumerate}
  \item \textbf{Revisión manual}: \menu{Respuestas por revisar} (panel) o la pestaña \menu{Revisión manual} de la
        sección. En cada respuesta puede poner un crédito manual (en \%) y un comentario.
  \item \textbf{Publicar notas}: en la configuración de la actividad marque \emph{Calificaciones publicadas} y guarde.
        Hasta entonces, el estudiante no ve la nota ni los comentarios.
  \item \textbf{Intentos abiertos tras el cierre}: en Resultados, \boton{Entregar y calificar} los califica con lo
        guardado a tiempo.
  \item \textbf{Exportar}: la pestaña Calificaciones exporta un CSV (separador «;» y coma decimal, para Excel en español).
\end{enumerate}

\chapter{Seguimiento de estudiantes}

\section{Pestaña Seguimiento}
Lista a quien puede necesitar apoyo, con el motivo y un botón para escribirle o abrir su ficha. Son avisos de apoyo
académico calculados con entregas y calificaciones, no indicadores de conducta.

\begin{center}
\begin{tabularx}{\linewidth}{>{\bfseries}l X}
\toprule
Aviso & Regla \\
\midrule
No entregó & Una actividad publicada de su grupo cerró sin entrega. \\
Cierra pronto & Una actividad cierra en menos de 48 horas y aún no la entregó. \\
Promedio bajo & Promedio por debajo del 70\,\%, con al menos dos actividades calificadas. \\
Notas bajando & Sus dos últimas notas están 20 puntos o más por debajo de las anteriores. \\
Intento abierto & Quedó un intento sin entregar al cerrar la actividad. \\
\bottomrule
\end{tabularx}
\end{center}

\section{Pendientes de entregar y recordatorios}
En los Resultados de cada actividad, \textbf{Pendientes de entregar} separa a quien no ha empezado de quien empezó sin
entregar. \boton{Enviar recordatorio} les manda un mensaje en vivo (texto editable, vigente 24 horas).

\section{Comunicación en vivo}
\ruta{Herramientas $\rightarrow$ Comunicación en vivo}: mensajes informativos, importantes o urgentes a una selección,
a una sección o a los conectados. Desde el teléfono, Orbe Control ofrece además el pulso de la clase y la pregunta
relámpago.

\chapter{Centro de Datos}
Desde \textbf{Orbe Control $\rightarrow$ Datos} se exportan los datos de una sección, un estudiante o una actividad en
CSV, JSON Lines, Excel o script SQLite, con un manifiesto (SHA-256), un README y un diccionario de datos.

\begin{itemize}
  \item \textbf{Modo identificado} o \textbf{modo investigación} (seudónimos generados en el servidor).
  \item Datos originales en lista blanca, incluidos los del motor de examen: configuración, sesiones, eventos
        oficiales y tramos de tiempo.
  \item Métricas derivadas, entre ellas \texttt{tiempos\_examen}: una fila por intento y pregunta con el tiempo de cada
        categoría, visitas, tiempo observable y prórroga.
  \item Cada exportación queda registrada.
\end{itemize}

\chapter{Laboratorio LIEIA}
Es el laboratorio de investigación sobre integridad académica. Sus indicadores son \textbf{experimentales}: describen
qué tan inusual sería un patrón si los estudiantes trabajaran de forma independiente, no prueban nada y nunca deciden
una calificación.

\begin{itemize}
  \item \textbf{D6 v1} (tiempo insensible a la dificultad) está \textbf{retirado} por defectos metodológicos; sus
        valores históricos aparecen en gris como «no usar».
  \item \textbf{D6 v2} es experimental y no se usa con datos reales hasta tener una referencia calibrada y una
        validación independiente.
  \item Los indicadores están congelados mientras se consolida la telemetría nueva.
\end{itemize}

% ============================================================================
\part{Administración}

\chapter{Puesta en producción}
Los cambios de base de datos y de funciones del servidor se publican con un procedimiento que protege los datos de
los estudiantes. Se ejecuta en la terminal, en la carpeta del proyecto:

\begin{enumerate}
  \item \texttt{npm run respaldo}: respaldo verificado de producción.
  \item \texttt{npm run probar}: aplica los cambios a una copia local con los datos reales, comprueba que no se
        pierde ninguna fila y ejecuta todas las pruebas de la base de datos. No se continúa si aparece algún error.
  \item \texttt{npm run plan}: muestra qué se aplicaría y un código de plan.
  \item \texttt{npm run aplicar -- --plan CÓDIGO}: nuevo respaldo, cada migración en su transacción y comprobación
        de los datos protegidos.
  \item \texttt{git push origin main}: publica la interfaz (GitHub Actions prueba, compila y verifica antes).
\end{enumerate}

\begin{importante}[Reglas que no se rompen]
\begin{itemize}
  \item Las migraciones son \textbf{aditivas}: nunca se borran tablas, columnas ni datos de estudiantes.
  \item Las credenciales y las claves de respuesta nunca van al repositorio ni a la interfaz.
  \item Ninguna calificación histórica se recalcula automáticamente.
\end{itemize}
\end{importante}

\section{Reversión}
\begin{itemize}
  \item \textbf{Motor de examen}: poner las evaluaciones en «Todas en una página»; si hiciera falta, una migración que
        restaure las reglas anteriores. Las tablas nuevas no se borran.
  \item \textbf{Prácticas asistidas}: dejar de crearlas; las tradicionales no cambian.
  \item \textbf{Interfaz}: revertir el commit y publicar.
\end{itemize}
El detalle técnico está en \texttt{docs/DESPLIEGUE\_REDISENO.md}.

% ============================================================================
\appendix

\chapter{Glosario}
\begin{description}[style=nextline]
  \item[Intento] Cada vez que un estudiante empieza una actividad.
  \item[Primer intento (práctica asistida)] La respuesta que el estudiante da antes de consultar ayudas; queda guardada
        sin cambios.
  \item[Rendimiento inicial declarado] Proporción de primeros intentos correctos, junto con lo que el estudiante
        declaró sobre su autonomía. No es una verificación.
  \item[Rendimiento posterior a la asistencia] Proporción de respuestas finales correctas después de estudiar con
        ayuda.
  \item[Transferencia] Ejercicio distinto con el mismo concepto, para comprobar comprensión.
  \item[Retención] Resultado en una comprobación posterior vinculada al tema.
  \item[Reloj oficial] La hora del servidor, que decide el límite de un examen.
  \item[Tramo de tiempo] Periodo del intento clasificado como trabajo, inactividad, pausa, página oculta, desconexión
        o desconocido.
  \item[Prórroga] Tiempo extra concedido por el profesor a un intento concreto.
  \item[Modo concentración] Durante un examen o el asistente, se oculta la navegación del aula en el teléfono.
\end{description}

\chapter{Preguntas frecuentes}
\begin{description}[style=nextline]
  \item[Se me fue la internet en un examen. ¿Pierdo lo que escribí?]
        No. Lo escrito se guarda en el dispositivo y se envía al volver la red. El reloj oficial sí sigue corriendo.
  \item[Cerré la página sin querer.]
        Vuelva a entrar en la evaluación: continúa en la misma pregunta con sus respuestas.
  \item[¿Me bajan la nota por decir que usé inteligencia artificial en una práctica asistida?]
        No. La declaración no interviene en la nota.
  \item[¿Puedo corregir mi primer intento en una práctica asistida?]
        El primer intento no cambia, pero en el paso «Aprender» puede dar una respuesta corregida, que es la que se
        califica.
  \item[Un estudiante tuvo un problema durante el examen. ¿Qué hago?]
        En \emph{En vivo}, pulse \boton{+5 min} en su fila y anote el motivo.
  \item[¿Por qué un estudiante aparece «Sin conexión»?]
        Lleva más de 45 segundos sin contacto con el servidor: red, pantalla apagada o página cerrada. Sus respuestas
        guardadas no se pierden.
\end{description}

\end{document}
